Nei capitoli precedenti è stato definito il quadro teorico di riferimento: dal modello temporale e le metriche di traffico (vedi Capitolo \ref{ch:ch4}), alla costruzione geometrica degli \emph{pseudo-settori collassati} (vedi Sezione \ref{subsec:algo-collassati}), fino alla formulazione del modello del \gls{dac} come modello di \gls{ilp} (vedi Sezione \ref{sec:dac_math}).

In questo capitolo si affronta la traduzione operativa di tali concetti, descrivendo l'architettura del prototipo software sviluppato per validare la metodologia proposta. L'attenzione si sposta quindi dal piano progettuale a quello realizzativo: l'obiettivo del capitolo non è fornire una documentazione esaustiva del codice sorgente, bensì motivare le scelte tecnologiche, illustrare l'integrazione delle librerie e descrivere l'organizzazione della pipeline di elaborazione dei dati.
Verranno dapprima presentati l'ambiente di sviluppo e gli strumenti adottati, per poi analizzare l'implementazione degli algoritmi di calcolo delle metriche e di generazione degli scenari, concludendo con una disamina delle prestazioni e della verifica funzionale.

\section{Ambiente di sviluppo e strumenti}
In questa sezione si descrive l'ambiente di sviluppo e l'insieme di strumenti software utilizzati per costruire la pipeline di elaborazione dei dati \gls{adrr} e il prototipo di risolutore per il problema di \gls{dac}.  

Le attività sono state svolte prevalentemente in ambiente \gls{python} su workstation personale. Tutti gli strumenti convergono in un unico progetto software organizzato come applicazione a riga di comando, per facilitare la riproducibilità degli esperimenti descritti nelle sezioni successive.

\subsection{Linguaggio di programmazione: Python}

La scelta del linguaggio di riferimento è ricaduta su \gls{python} (versione 3.13), un linguaggio di programmazione di alto livello, interpretato e orientato agli oggetti \cite{python} (Figura~\ref{fig:python_logo}).

\begin{figure}[htbp]
    \centering
    \includegraphics[width=0.2\textwidth]{chapters/ch_impl/img/python.pdf} 
    \caption{Logo di Python.}
    \label{fig:python_logo}
\end{figure}

\gls{python} si distingue per la sua sintassi chiara e leggibile, ma soprattutto per un vastissimo ecosistema di librerie di terze parti che lo rendono uno standard per la data science.
La decisione di utilizzare \gls{python}, preferendolo ad alternative pur validissime in ambito scientifico come R \cite{r_language} o MATLAB \cite{matlab}, è motivata dalla necessità di integrare in un unico stack tecnologico domini molto diversi: l'elaborazione di grandi moli di dati, la manipolazione geospaziale e l'interfacciamento con solutori di ottimizzazione.

L'intero progetto è stato organizzato come applicazione a riga di comando (\acrshort{cli}), eseguita in ambienti virtuali isolati per garantire la gestione delle dipendenze e la riproducibilità degli esperimenti.

\subsection{Librerie per l'elaborazione dei dati}

I rilasci \gls{adrr} sono disponibili in formato \acrshort{csv} compresso, che costituisce il formato di input obbligato. Per i risultati intermedi sono stati invece adottati file \texttt{Parquet}, che offrono una lettura più efficiente, una migliore compressione e un’integrazione diretta con \texttt{pandas} e \texttt{Dask}. Non è stato introdotto un \gls {dbms} relazionale: l’onere di modellazione e caricamento non sarebbe stato compensato, nel contesto dello stage, da vantaggi concreti rispetto a una pipeline basata su file e prevalentemente \emph{read–only}.
La gestione dei dati \gls{adrr}, che costituiscono l'input del sistema, richiede strumenti capaci di manipolare strutture tabellari complesse e di gestire volumi di dati che possono superare la memoria disponibile su una singola workstation (Figura~\ref{fig:data_libs}).

\subsubsection{Pandas}
\texttt{pandas} (Figura~\ref{fig:pandas_logo}) è una libreria \gls{open source} che fornisce strutture dati ad alte prestazioni e strumenti di analisi facili da utilizzare \cite{pandas}. La struttura fondamentale introdotta è il \texttt{DataFrame}, una tabella bidimensionale con etichette per righe e colonne, simile a un foglio di calcolo o a una tabella \acrshort{sql}, ma manipolabile in memoria.
Nel contesto del progetto, \texttt{pandas} è lo strumento primario per il filtraggio, l'aggregazione e il \emph{joining} dei dati relativi ai profili di volo e alle metriche di traffico su finestre temporali limitate.

\subsubsection{Dask}
\texttt{Dask} (Figura~\ref{fig:dask_logo}) è una libreria per il calcolo parallelo flessibile in \gls{python} \cite{dask}. Essa estende le funzionalità di \texttt{pandas} permettendo di lavorare su dataset "\emph{larger-than-memory}" (più grandi della \acrshort{ram} disponibile). \texttt{Dask} suddivide i grandi dataset in molteplici piccole partizioni, gestendone l'esecuzione in parallelo e in modalità \emph{lazy} (i calcoli vengono eseguiti solo quando il risultato è effettivamente richiesto).
L'adozione di \texttt{Dask} è risultata essenziale per i passaggi più onerosi della pipeline, come il caricamento e l'elaborazione di interi mesi di traffico aereo \gls{adrr}.

\begin{figure}[htbp]
    \centering
    \begin{subfigure}{0.45\textwidth}
        \centering
        \includegraphics[height=2.5cm]{chapters/ch_impl/img/pandas.pdf} % Logo Pandas
        \caption{Pandas}
        \label{fig:pandas_logo}
    \end{subfigure}
    \hfill
    \begin{subfigure}{0.45\textwidth}
        \centering
        \includegraphics[height=2.5cm]{chapters/ch_impl/img/dask.pdf} % Logo Dask
        \caption{Dask}
        \label{fig:dask_logo}
    \end{subfigure}
    \caption{Loghi delle librerie utilizzate per la manipolazione dei dati tabellari.}
    \label{fig:data_libs}
\end{figure}


\subsection{Librerie Geospaziali}
La costruzione delle porzioni di spazio aereo richiede la gestione di geometrie complesse e relazioni topologiche (intersezioni, contenimento, adiacenza).

\subsubsection{Shapely}
\texttt{Shapely} è una libreria per la manipolazione e l'analisi di oggetti geometrici planari \cite{shapely}. Fornisce le primitive geometriche fondamentali (punti, linee, poligoni) e i metodi per operare su di esse (come \texttt{union}, \texttt{intersection}, \texttt{difference}).
Nel progetto, \texttt{Shapely} è il motore geometrico utilizzato per calcolare fisicamente le intersezioni tra le traiettorie di volo e i confini dei settori.

\subsubsection{GeoPandas}
\texttt{GeoPandas} è un progetto che estende i tipi di dati di \texttt{pandas} per consentire operazioni spaziali su oggetti geometrici \cite{geopandas} (Figura~\ref{fig:cplex_logo}). Integra \texttt{Shapely} per le operazioni geometriche, consentendo di associare a ogni riga di un \texttt{DataFrame} una geometria.
L'uso di \texttt{GeoPandas} ha permesso di trattare le regioni dello spazio aereo come \texttt{GeoDataFrame}, semplificando l'integrazione con i dati di traffico.

\begin{figure}[htbp]
    \centering
    % File geopandas_logo.svg
    \includegraphics[width=0.45\textwidth]{chapters/ch_impl/img/geopandas.pdf} 
    \caption{Logo di GeoPandas.}
    \label{fig:geo_logo}
\end{figure}

\subsection{Ottimizzazione matematica}
Lo scopo finale dei dati elaborati consiste nella risoluzione del problema \gls{dac}, modellato come problema di \gls{ilp}.

\subsubsection{IBM ILOG CPLEX}
\gls{cplex} è uno dei solutori di ottimizzazione più potenti e diffusi in ambito accademico e industriale \cite{cplex} (Figura~\ref{fig:cplex_logo}). È progettato per risolvere problemi di programmazione lineare e mista-intera di grandi dimensioni con elevata efficienza. È stato preferito a soluzioni \gls{open source} per le sue prestazioni superiori su istanze complesse e per la disponibilità di una licenza accademica.

\subsubsection{DOcplex}
Per interfacciare il modello matematico con il motore di risoluzione, è stata utilizzata \texttt{DOcplex}, la libreria nativa di \acrshort{ibm} \cite{docplex}. Questa permette di definire variabili e vincoli utilizzando una sintassi \gls{python} naturale che rispecchia la notazione matematica, traducendo poi il modello per il motore \gls{cplex}.

\begin{figure}[htbp]
    \centering
    % File cplex_logo.svg
     \includegraphics[width=0.35\textwidth]{chapters/ch_impl/img/ilog.pdf} 
    \caption{Logo di IBM ILOG CPLEX.}
    \label{fig:cplex_logo}
\end{figure}

\subsection{Controllo di versione: Git}

La gestione del ciclo di vita del software è stata affidata a \gls{git} (Figura~\ref{fig:git_logo}), il sistema di controllo di versione distribuito standard per lo sviluppo moderno \cite{git}.

\begin{figure}[htbp]
    \centering
    % Assicurati che il file si chiami git_logo.svg
    \includegraphics[width=0.35\textwidth]{chapters/ch_impl/img/git.pdf} 
    \caption{Logo ufficiale di Git.}
    \label{fig:git_logo}
\end{figure}

L'utilizzo di \gls{git} ha garantito non solo il salvataggio sicuro dello storico delle modifiche, ma ha anche abilitato un flusso di lavoro basato su \emph{branching} (ramificazioni). Questo approccio ha permesso di sviluppare e testare nuove euristiche di ottimizzazione o refactoring strutturali in ambienti isolati (branch), per poi integrarle nel ramo principale (\emph{main}) solo una volta verificate stabilità e correttezza. Tale metodologia è risultata cruciale per mantenere l'integrità della base di codice durante le fasi esplorative.

\subsection{Redazione del documento: \LaTeX}

La stesura del presente elaborato è stata realizzata mediante \LaTeX (Figura~\ref{fig:latex_logo}), un sistema di preparazione di testi basato sul programma di composizione tipografica \TeX\ \cite{latex_project}.

\begin{figure}[htbp]
    \centering
    % Assicurati che il file si chiami latex_logo.svg
    \includegraphics[width=0.35\textwidth]{chapters/ch_impl/img/latex.pdf} 
    \caption{Logo del progetto \LaTeX.}
    \label{fig:latex_logo}
\end{figure}

A differenza dei tradizionali elaboratori di testo \gls{wysiwyg}, \LaTeX\ adotta un approccio basato sulla separazione tra contenuto e stile. Questa caratteristica lo rende lo strumento ideale per la redazione scientifica, garantendo una gestione automatizzata e coerente di riferimenti incrociati, bibliografia (tramite Bib\TeX), glossario e formule matematiche complesse.
Per questo progetto è stata utilizzata una classe personalizzata per garantire la conformità agli standard dell'ateneo e l'accessibilità del documento finale (PDF/A).

\section{Implementazione}
In questa sezione si descrive come siano stati tradotti in codice \gls{python} gli schemi algoritmici introdotti nei Capitoli \ref{ch:ch4} e \ref{ch:ch5} per il calcolo del \emph{throughput} (vedi Sezione \ref{subsec:throughput-teoria}), dell'\emph{occupancy} (vedi Sezione \ref{subsec:occupancy-teoria}), delle capacità nominali dei pseudo-settori, le procedure per la generazione delle configurazioni di \emph{pseudo-settori collassati}, e l’implementazione del modello di \gls{ilp} per il \gls{dac}.

\subsection{Throughput}
La funzione \texttt{compute}, appartenente al modulo \texttt{throughput}, orchestra l'implementazione dell'Algoritmo~\ref{alg:throughput-anyW-naturale}. Essa traduce i parametri matematici in oggetti \gls{python}: l'orizzonte temporale $[\tau_0, \tau_1]$, il parametro di discretizzazione $\lambda$ e la finestra $W$ sono mappati rispettivamente sugli argomenti \texttt{tau0}, \texttt{tau1}, \texttt{lambda} e \texttt{W}. Il DataFrame \texttt{E} rappresenta l'insieme degli eventi in ingresso.

La firma di seguito riportata evidenzia come la logica sia incapsulata per restituire un DataFrame strutturato contenente il calcolo finale:

\begin{lstlisting}[language=Python, caption={}]
def compute(
    E: dd.DataFrame,
    *,
    R_col: str,
    tau0: pd.Timestamp,
    tau1: pd.Timestamp,
    lambda: pd.Timedelta,
    W: pd.Timedelta,
    t_in_col: str = "Entry Time",
    expected_keys: Optional[List[str]] = None,
) -> pd.DataFrame:
\end{lstlisting}

La prima fase operativa consiste nella costruzione della griglia temporale. Attraverso la funzione ausiliaria \texttt{grid\_params}, vengono derivati i limiti operativi \texttt{t0} e \texttt{t1} e il numero di intervalli \texttt{K}. È fondamentale notare come l'orizzonte venga esteso fino a $\tau_1 + W$ (gestito internamente da \texttt{W\_td}) per garantire la corretta valutazione delle finestre che si sovrappongono al termine dell'osservazione. Sulla base di questi estremi, si procede alla creazione di \texttt{E\_filtered}, un sottoinsieme del DataFrame originale che trattiene esclusivamente gli eventi utili al calcolo corrente, riducendo il carico di memoria. Il seguente listato illustra la procedura di inizializzazione e filtraggio:

\begin{lstlisting}[language=Python, caption={}]
t0, t1, Lambda_td, W_td, K, _m = grid_params(
    tau0, tau1, lambda, str(W)
)

E_filtered = E[[R_col, t_in_col]]
E_filtered = E_filtered[
    (E_filtered[t_in_col] >= t0) &
    (E_filtered[t_in_col] <  t1 + W_td)
]
\end{lstlisting}

Una volta isolati i dati pertinenti, il flusso si sposta sull'aggregazione. Poiché l'algoritmo richiede la valutazione della funzione di conteggio $N_R(t)$, è necessario prima ridurre gli eventi a una serie temporale di conteggi discreti. Utilizzando un'operazione di \texttt{groupby} sulle colonne chiave e temporali (rispettivamente \texttt{R\_col}, \texttt{t\_in\_col}), si ottiene la variabile \texttt{events\_counted}, che contiene la cardinalità $n_R$ per ogni istante univoco.

A questo punto, sono necessarie due operazioni preliminari per garantire la coerenza matematica del calcolo successivo. In primo luogo, la colonna identificativa \texttt{R\_col} viene normalizzata al tipo \texttt{object} per prevenire conflitti durante le operazioni di unione. In secondo luogo, il DataFrame viene esplicitamente ordinato temporalmente tramite \texttt{sort\_values}. Questo ordinamento è un prerequisito stringente: affinché la somma cumulativa (\texttt{cumsum}) abbia senso fisico e affinché la successiva unione asincrona funzioni, la sequenza degli eventi deve essere rigorosamente cronologica. Solo a valle di questa preparazione viene calcolata la colonna \texttt{N\_R}, che rappresenta la cumulata degli eventi osservati fino a quel momento. Il codice che realizza tale aggregazione e il successivo conteggio è riportato di seguito:

\begin{lstlisting}[language=Python, caption={}]
events_counted = (
    E_filtered.groupby([R_col, t_in_col], observed=False)
        .size()
        .rename("n_R")
        .reset_index()
        .compute()  # Esecuzione del calcolo
)
events_counted[R_col] = events_counted[R_col].astype("object")
events_counted = events_counted.sort_values(t_in_col).reset_index(drop=True)
events_counted["N_R"] = events_counted.groupby(
    R_col, sort=False, observed=False
)["n_R"].cumsum()
\end{lstlisting}

Il cuore della funzione risiede nella necessità di riconciliare due assi temporali disallineati:
\begin{itemize}
    \item La griglia di misurazione, definita dagli istanti artificiali ed equispaziati $t_k$
    \item La storia degli eventi, costituita dagli istanti stocastici (casuali) $t_{in}$ in cui avvengono le osservazioni reali
\end{itemize}
Teoricamente, la grandezza cumulativa $N_R(t)$ è modellabile come una funzione a gradini (\emph{step function}) continua a tratti: il suo valore, che rappresenta il conteggio totale degli eventi, rimane costante (fase di \emph{hold}) tra un arrivo e l'altro, incrementandosi solo al verificarsi di un nuovo evento.

Per calcolare il throughput $T_R(k)$, è necessario campionare il valore di questa cumulata esattamente in corrispondenza dei "tagli" imposti dalla griglia ($t_k$ e $t_k+W$). Tuttavia, è statisticamente improbabile che un evento reale coincida al microsecondo con un istante della griglia; di conseguenza, un'operazione di \textit{join} tradizionale (basata sull'uguaglianza esatta delle chiavi temporali) fallirebbe, generando valori mancanti.

L'implementazione risolve questa criticità utilizzando la funzione \texttt{pd.merge\_asof}. Essa opera come segue: per ogni istante $t$ della griglia, il sistema guarda all'indietro (\texttt{direction="backward"}) e recupera l'ultimo valore valido di $N_R$ registrato prima di $t$. 

Inoltre, impostando \texttt{allow\_exact\_matches=False}, si forza l'algoritmo a escludere l'istante attuale, soddisfacendo così la definizione formale di cumulata stretta ($< t$).

Questa logica viene applicata in due fasi distinte. Per il bordo sinistro dell'intervallo si utilizza direttamente la colonna \texttt{t\_k}; per il bordo destro, viene calcolata una colonna temporanea traslata (\texttt{t\_k\_plus\_W}) su cui effettuare l'allineamento. Le istruzioni riportate di seguito implementano l'ordinamento preliminare e la definizione della struttura dati per per il calcolo ai bordi:
\begin{lstlisting}[language=Python, caption={}]
# Ordinamento preliminare degli eventi (necessario per merge_asof)
events_sorted = events_counted.sort_values(t_in_col).reset_index(drop=True)

# Valutazione di N_R(t_k)
# merge_asof tra la griglia (t_k) e gli eventi (t_in_col)
left_df = base_grid[[R_col, "t_k", "idx"]].sort_values("t_k")

left_merged = pd.merge_asof(
    left_df, events_sorted,
    left_on="t_k", right_on=t_in_col,
    by=R_col, 
    direction="backward",       # Cerca nel passato
    allow_exact_matches=False,  # Implementa la disuguaglianza stretta (< t)
).sort_values("idx").reset_index(drop=True)

# Valutazione di N_R(t_k + W)
# Si crea un riferimento temporale traslato: t_k + W
right_df = base_grid[[R_col, "t_k", "idx"]].copy()
right_df["t_k_plus_W"] = right_df["t_k"] + W_td
right_df = right_df.sort_values("t_k_plus_W")

right_merged = pd.merge_asof(
    right_df, events_sorted,
    left_on="t_k_plus_W", right_on=t_in_col,
    by=R_col, 
    direction="backward",
    allow_exact_matches=False,
).sort_values("idx").reset_index(drop=True)
\end{lstlisting}

Con i valori delle cumulate correttamente campionati sulla griglia, il calcolo si riduce a una differenza vettoriale tra le colonne estratte dai due DataFrame uniti. I vettori risultanti vengono sottratti per ottenere \texttt{T\_R}, gestendo eventuali valori mancanti tramite riempimento con zeri.

Infine, la funzione assembla il DataFrame di output \texttt{T\_output}, associando a ogni riga i riferimenti temporali espliciti e il valore calcolato. Di seguito viene riportata l'implementazione del calcolo differenziale e la preparazione dell'output:

\begin{lstlisting}[language=Python, caption={}]
# T_R(k) = N_R(t_k+W) - N_R(t_k)
N_R_at_tk        = left_merged["N_R"].fillna(0).to_numpy(dtype=np.int64)
N_R_at_tk_plus_W = right_merged["N_R"].fillna(0).to_numpy(dtype=np.int64)
T_R = N_R_at_tk_plus_W - N_R_at_tk

T_output = base_grid.drop(columns=["idx"]).copy()
T_output["t_k_plus_W"] = T_output["t_k"] + W_td
T_output["T_R"] = T_R
return T_output[[R_col, "t_k", "t_k_plus_W", "T_R"]]
\end{lstlisting}


\subsection{Occupancy}
Analogamente a quanto visto per il \emph{throughput}, la funzione \texttt{compute} del modulo \texttt{occupancy} traduce l'Algoritmo~\ref{alg:occupancy-anyW-diff} in istruzioni vettoriali ottimizzate per \texttt{Dask} e \texttt{Pandas}.
L'obiettivo è calcolare l'occupazione $O_R(k)$, intesa come il numero di velivoli che risultano presenti in \(R\) per tutta la finestra $[t_k, t_k + W)$:

La firma della funzione di seguito riportata riflette la necessità di mappare i concetti teorici (l'insieme degli eventi $E$, il parametro di discretizzazione $\lambda$ e la finestra $W$) sui tipi di dati \gls{python}:
\clearpage

\begin{lstlisting}[language=Python, caption={}]
def compute(
    E: dd.DataFrame,
    *,
    key_col: str,
    tau0: pd.Timestamp,
    tau1: pd.Timestamp,
    Delta: pd.Timedelta,
    W: pd.Timedelta,
    t_in: str = "Entry Time",
    t_out: str = "Exit Time",
) -> pd.DataFrame:
\end{lstlisting}

Dopo la fase preliminare di definizione della griglia temporale (identica a quella descritta per il \emph{throughput}), l'algoritmo adotta un approccio differenziale per massimizzare l'efficienza. Invece di verificare per ogni istante $t_k$ quali eventi siano attivi (operazione che richiederebbe complessità quadratica o join costosi), ogni evento $i$ contribuisce al carico del sistema a partire da un indice di inizio $k_{start}$ fino a un indice di fine $k_{end}$.

Il codice calcola questi indici discretizzando i tempi di ingresso ($t_{in}$) e uscita ($t_{out}$) rispetto all'origine $\tau_0$ e al passo $\lambda$. L'uso delle funzioni \texttt{ceil} per l'ingresso e \texttt{floor} per l'uscita garantisce il rispetto rigoroso delle condizioni al contorno della finestra semi-aperta. L'implementazione è riportata di seguito:

\begin{lstlisting}[language=Python, caption={}]
# Calcolo vettoriale degli indici sulla griglia
k_start = ceil((t_in - t0)/lambda)

k_end = floor((t_out - W - t0)/lambda)

\end{lstlisting}

Una volta definiti gli intervalli di validità per ciascun evento, il problema viene trasformato in una somma di impulsi. Si costruiscono due DataFrame temporanei: il primo, \texttt{e1}, rappresenta gli "ingressi" e assegna un valore $+1$ all'indice $k_{start}$; il secondo, \texttt{e2}, rappresenta le "uscite" e assegna un valore $-1$ all'indice immediatamente successivo alla fine dell'evento ($k_{end} + 1$).

Questa tecnica riduce l'intero dataset a una serie di variazioni locali ($\delta$). Concatenando questi contributi, si ottiene un flusso di variazioni che descrive come l'occupazione cambia nel tempo. L'implementazione operativa di tale meccanismo a impulsi è riportata nel frammento sottostante:
\clearpage
\begin{lstlisting}[language=Python, caption={}]
e1 = dfv[[key_col, "k_start"]].rename(columns={"k_start": "k"})
e1["delta"] = 1

# L'uscita avviene dopo k_end, quindi il -1 si applica a k_end + 1
e2 = dfv[[key_col, "k_end"]].rename(columns={"k_end": "k"})
e2["k"] = e2["k"] + 1
e2["delta"] = -1

edges_dd = dd.concat([e1, e2], axis=0)
# Filtro per mantenere solo gli indici all'interno dell'orizzonte [0, K)
edges_dd = edges_dd[(edges_dd["k"] >= 0) & (edges_dd["k"] < K)]
\end{lstlisting}

L'ultimo passaggio consiste nella ricostruzione del segnale di occupazione. Le variazioni vengono raggruppate per regione e per istante temporale, sommate algebricamente, e infine integrate tramite una somma cumulativa (\texttt{cumsum}).

Per garantire che il risultato copra l'intera griglia temporale (anche laddove non ci sono variazioni), le somme parziali vengono unite (\textit{merge}) a una struttura base completa ("scheletro") che contiene tutte le combinazioni possibili di regioni $R$ e istanti $t_k$. La colonna risultante \texttt{O\_R} rappresenta l'integrale discreto delle variazioni, ovvero il numero esatto di eventi presenti nel sistema in ogni finestra temporale. Il listato conclusivo illustra la procedura di integrazione tramite somma cumulativa e la composizione dell'output finale:
\begin{lstlisting}[language=Python, caption={}]
# Aggregazione delle variazioni e unione con la griglia base
base = base.merge(edges, how="left", on=[key_col, "k"]).fillna({"delta": 0})
base = base.sort_values([key_col, "k"]).reset_index(drop=True)

# Calcolo finale tramite somma cumulativa
base["O_R"] = base.groupby(key_col, sort=False)["delta"].cumsum().astype("int64")

# Associazione dei timestamp espliciti
base["t_k"] = np.tile(tks.to_numpy(), len(Rs))
base["t_k_plus_W"] = base["t_k"] + W
return base[[key_col, "t_k", "t_k_plus_W", "O_R"]]
\end{lstlisting}


\clearpage
\subsection{Calcolo delle capacità}
L'implementazione delle capacità nominali traduce in codice i passi definiti dagli Algoritmi~\ref{alg:cap-elementare-naturale} e~\ref{alg:cap-collassato-naturale}. Il processo si articola in due fasi: la stima delle capacità per i \emph{pseudo-settori elementari}, basata sull'analisi storica del throughput, e la successiva derivazione delle capacità per i \emph{pseudo settori collassati}.

La procedura inizia con l'ingestione dei dati di \emph{throughput} calcolati secondo la logica dell'Algoritmo \ref{alg:throughput-anyW-naturale}. Poiché la capacità è definita su base giornaliera, il dataset viene filtrato per isolare una specifica finestra temporale di 24 ore (variabili \texttt{day\_start} e \texttt{day\_end}), avendo cura di normalizzare tutti i riferimenti temporali al fuso orario UTC per garantire la coerenza con i dati di volo. Il seguente frammento illustra le operazioni di caricamento e selezione temporale dei dati:
\begin{lstlisting}[language=Python, caption={}]
df_throughput = pd.read_parquet(throughput_parquet_path)
df_throughput["t_k"] = pd.to_datetime(df_throughput["t_k"], utc=True)

day_start = pd.Timestamp(day, tz="UTC")
day_end = day_start + pd.Timedelta(days=1)
df_day = df_throughput[
    (df_throughput["t_k"] >= day_start) & 
    (df_throughput["t_k"] < day_end)
]
\end{lstlisting}
Per determinare la capacità di un \emph{pseudo-settore elementare}, non è sufficiente osservare il picco massimo assoluto (spesso un outlier) né la media semplice. L'algoritmo procede iterando su ogni regione $R$ e applicando filtri statistici.

Inizialmente, vengono scartati gli intervalli a traffico nullo. Successivamente, per rimuovere sia il rumore di fondo che i picchi anomali non sostenibili, si considerano solo i valori di throughput compresi tra il 10° e il 90° percentile (\texttt{p\_10} e \texttt{p\_90}). Un ulteriore filtro scarta i valori inferiori a una soglia minima (\texttt{threshold}), utile per evitare stime falsate in settori scarsamente utilizzati.

La capacità viene stimata calcolando la media del quartile superiore (il top 25\% dei valori ordinati). Il valore finale include un termine di aggiustamento aleatorio (\texttt{random\_adjustment}) per introdurre variabilità realistica nelle simulazioni, con un meccanismo di fallback (distribuzione uniforme) qualora i dati storici fossero insufficienti. L'implementazione iterativa di questa procedura di calcolo è mostrata nel seguente listato:
\clearpage
\begin{lstlisting}[language=Python, caption={}]
capacities = {}
for region_id, group in df_day.groupby("R"):
    counts = group["T_R"]
    counts = counts[counts != 0]  # Rimozione periodi inattivi
    
    # Gestione settori vuoti
    if counts.empty:
        capacities[region_id] = npr.uniform(threshold * 0.5, threshold * 1.5)
        continue
    
    # Filtro statistico (intervallo interquartile espanso)
    p10 = counts.quantile(0.10)
    p90 = counts.quantile(0.90)
    filtered_counts = counts[(counts >= p10) & (counts <= p90)]
    
    # Rimozione traffico non significativo
    filtered_counts = filtered_counts[filtered_counts >= threshold]
    if filtered_counts.empty:
        capacities[region_id] = npr.uniform(threshold * 0.5, threshold * 1.5)
        continue

    # Calcolo sul quartile superiore (High Sustained Throughput)
    sorted_counts = filtered_counts.sort_values()
    num_top = max(1, int(round(len(sorted_counts) * 0.75)))
    top_counts = sorted_counts.iloc[-num_top:]
    avg_top = top_counts.mean()
    
    capacity = avg_top + random_adjustment
	capacities[region_id] = (
    	capacity if capacity >= 1 
    	else npr.uniform(threshold * 0.5, threshold * 1.5)
    )
\end{lstlisting}

Una volta calcolate le capacità dei \emph{pseudo-settori elementari}, il codice calcola quelle dei \emph{pseudo-settori collassati}. In questo contesto, la capacità non è la semplice somma delle parti (che risulterebbe irrealisticamente alta).

Per ogni \emph{pseudo-settore collassato}, si recuperano le capacità dei suoi costituenti. La capacità risultante è definita partendo dal valore massimo tra questi (\texttt{b\_capacity}), incrementato di un fattore proporzionale al numero di settori aggiuntivi (10\% per ogni settore extra, fino a un massimo del 50\%). Questa euristica riflette il guadagno ottenuto gestendo un volume di spazio aereo più ampio, pur riconoscendo i strutturali del controllo del traffico aereo:

La capacità base è definita come media del 25\% più elevato dei valori superstiti. L'aggiustamento aleatorio finale introduce variabilità controllata, con ulteriore fallback per valori non validi.
Per i \emph{pseudo-settori collassati}, la capacità deriva da quelle elementari costituenti. Le seguenti istruzioni calcolano la capacità aggregata applicando il fattore di incremento proporzionale:

\begin{lstlisting}[language=Python, caption={}]
global_capacities = {}
for key, info in global_collapsed_info.items():
    elementary_ids = info["elementary_ids"]
    
    capacities_list = [element_capacities[es_id] 
                       for es_id in elementary_ids 
                       if es_id in element_capacities]

    if capacities_list:
        b_capacity = max(capacities_list)
        
        # Euristiche di aggregazione
        additional_count = len(elementary_ids) - 1
        increment_factor = min(additional_count * 0.1, 0.5)
        
        collapsed_capacity = b_capacity * (1 + increment_factor)
    else:
        # Fallback per settori senza dati
        collapsed_capacity = npr.uniform(25.0 * 0.5, 25.0 * 1.5)
    
    global_capacities[info["collapsed_sector"]] = int(round(collapsed_capacity))
\end{lstlisting}
\subsection{Generazione configurazioni di pseudo-settori collassati}
\label{sec:gen}
La procedura di generazione delle configurazioni collassate traduce l'Algoritmo~\ref{alg:collasso-dfs-prob} in una pipeline di elaborazione spaziale. Il processo inizia con la ricostruzione della geometria dei pseudo-settori elementari a partire dai dati grezzi.

La funzione \texttt{load\_fir\_data} legge il file CSV, raggruppa le coordinate per identificativo (\texttt{Airspace ID}) e le converte in poligoni chiusi. Durante questa fase, vengono estratti gli attributi verticali (\gls{fl} minimi e massimi) e le geometrie vengono proiettate in un sistema di riferimento metrico (EPSG:3035) per consentire calcoli di adiacenza precisi. L'uso di \texttt{unary\_union} assicura che eventuali geometrie frammentate vengano consolidate in un unico oggetto spaziale valido. Il codice seguente implementa la funzione di importazione e normalizzazione geometrica appena descritta:
\clearpage
\begin{lstlisting}[language=Python, caption={}]
def load_fir_data() -> gpd.GeoDataFrame:
    df = pd.read_csv(DEFAULT_FIR_FILE, sep=",", compression="gzip",
                     quotechar='"', skipinitialspace=True)

    for fir_id, group in df.groupby("Airspace ID"):
        group = group.sort_values("Sequence Number")  
        coords = list(zip(group["Longitude"], group["Latitude"]))
        
        if coords[0] != coords[-1]:
            coords.append(coords[0])
               
        line = LineString(coords)
        polys = list(polygonize([line]))
        geom_union = unary_union([p.buffer(0) for p in polys]) if polys else None
        
        records.append({
            "Airspace ID": fir_id,
            "Min Flight Level": group["Min Flight Level"].iloc[0],
            "Max Flight Level": group["Max Flight Level"].iloc[0],
            "geometry": geom_union
        })

    return gpd.GeoDataFrame(records, crs="EPSG:4326").to_crs("EPSG:3035")
\end{lstlisting}
Prima di procedere all'aggregazione, è necessario definire rigorosamente il criterio di adiacenza. Poiché i settori sono volumi tridimensionali, due regioni adiacenti sulla mappa possono essere considerate "vicine" solo se condividono anche una porzione significativa di spazio verticale. La funzione \texttt{vertical\_overlap} implementa questo vincolo calcolando l'intersezione lineare tra gli intervalli di quota; solo se questa intersezione supera una soglia operativa definita (\texttt{VERTICAL\_OVERLAP\_MARGIN}), i settori sono considerati candidati per la fusione. Di seguito è riportata la codifica di tale verifica:
\begin{lstlisting}[language=Python, caption={}]
def vertical_overlap(row1: pd.Series, row2: pd.Series) -> bool:
    min1, max1 = row1["Min Flight Level"], row1["Max Flight Level"]
    min2, max2 = row2["Min Flight Level"], row2["Max Flight Level"]
    overlap = min(max1, max2) - max(min1, min2)
    return overlap >= VERTICAL_OVERLAP_MARGIN  
\end{lstlisting}
L'algoritmo di generazione è implementato nel metodo \texttt{generate}. L'approccio è di tipo stocastico-espansivo: si itera sui settori non ancora assegnati utilizzandoli come "semi" per nuovi cluster. Per ottimizzare la ricerca dei vicini, si sfrutta un indice spaziale (\texttt{sindex}), che riduce il costo computazionale delle query geometriche (\texttt{predicate="touches"}).

Una volta identificati i vicini spazialmente e verticalmente compatibili, l'algoritmo decide se espandere il cluster basandosi su una probabilità dinamica \texttt{p\_merge}. Questa probabilità è inversamente proporzionale alla dimensione attuale del gruppo: clusters piccoli tendono a crescere facilmente, mentre clusters vicini alla dimensione massima (\texttt{max\_size}) hanno una probabilità di espansione che tende a zero.

Al termine del ciclo di espansione, le geometrie dei settori costituenti vengono fuse tramite \texttt{unary\_union} per creare la forma finale del \emph{pseudo-settore collassato}. Il listato seguente mostra l'implementazione completa del ciclo di espansione e fusione:
\begin{lstlisting}[language=Python, caption={}]
def generate(self) -> gpd.GeoDataFrame:
    merged_indices = set() 
    merged_records = []
    sindex = self.sectors.sindex 
       
    for idx, row in self.sectors.iterrows():  
        if idx in merged_indices: continue              
        collapsed_sector = {idx}  
        stack = [idx]             
        
        # Espansione ricorsiva del gruppo
        while stack and len(collapsed_sector) < self.max_size:
            current_idx = stack.pop()  # Estrazione settore corrente dallo stack
            
            # Ricerca dei vicini spaziali tramite indice
            possible_idxs = list(sindex.query(
                self.sectors.loc[current_idx, "geometry"], 
                predicate="touches"
            ))
            
            # Valutazione dei vicini candidati
            for neighbor in possible_idxs:  
                if neighbor in merged_indices: continue  
                                    
                if not self.vertical_overlap(
                    self.sectors.loc[current_idx],
                    self.sectors.loc[neighbor]
                ): continue
                
                group_size = len(collapsed_sector)
                p_merge = (self.max_size - group_size + 1) / \
                          (self.max_size + 1)
                
                # Decisione probabilistica di inclusione
                if npr.rand() < p_merge:  
                    collapsed_sector.add(neighbor)  
                    stack.append(neighbor)         
                    if len(collapsed_sector) >= self.max_size: 
                        break  
        
        # Unione geometrica dei settori del gruppo
        merged_geom = unary_union([
            self.sectors.loc[i, "geometry"].buffer(0) 
            for i in collapsed_sector
        ])
        
        merged_records.append({
            "elementary_sector": list(collapsed_sector),
            "geometry": merged_geom,
            "collapsed_sector": generate_id()  # ID univoco
        })
        merged_indices.update(collapsed_sector)  
            
    return gpd.GeoDataFrame(merged_records, crs=self.sectors.crs)
\end{lstlisting}

\subsection{Implementazione del modello di ILP per il DAC}
L'implementazione del modello per il \gls{dac} traduce la formulazione matematica presentata nella Sezione \ref{sec:dac_math} in un linguaggio di programmazione, sfruttando la libreria \texttt{DOcplex} di \acrshort{ibm}. 

Il flusso di lavoro inizia con l'ingestione dei dati operativi. Per ogni scenario di simulazione, identificato da un \texttt{config\_id}, il sistema carica la domanda di traffico, le capacità stimate dei pseudo-settori e il grafo delle transizioni ammissibili. Questi dati grezzi vengono processati per quantificare il sovraccarico locale (eccesso di domanda rispetto alla capacità) e aggregarlo a livello di configurazione, ottenendo i costi che guideranno l'ottimizzazione.

Segue l'implementazione delle funzioni per il calcolo dell'eccesso di domanda e dei coefficienti di costo:
\clearpage
\begin{lstlisting}[language=Python, caption={}]
tr, c, dizio_conf = load_data_for_config(config_id, data_dir)

def compute_excess(tr, c, T):
    excess = {}
    for s in tr:
        excess[s] = [max(tr[s][t] - c[s], 0) for t in T]
    return excess

def compute_coefficients(dizio_conf, excess, T):
    e = {}
    for conf, sectors in dizio_conf.items():
        e[conf] = {t: sum(excess[s][t] for s in sectors) for t in T}
    return e
\end{lstlisting}

Definiti i parametri di costo, si procede alla costruzione dello spazio delle soluzioni istanziando le variabili decisionali. Il modello richiede due set di variabili: $x$, che indica lo stato attivo della configurazione, e $s$, variabile ausiliaria che traccia i cambiamenti di stato (switch). Sebbene il problema sia intrinsecamente binario (\gls{ilp}), l'implementazione offre la flessibilità di rilassare il dominio al continuo.

Di seguito è riportata la funzione per l'inizializzazione delle variabili nel modello:

\begin{lstlisting}[language=Python, caption={}]
def create_decision_variables(model, X, dom):
    if dom == 'cont':
        x = model.continuous_var_dict(keys=X, name='x', lb=0, ub=1)
        s = model.continuous_var_dict(keys=X, name='s', lb=0, ub=1)
    elif dom == 'bin':
        x = model.binary_var_dict(keys=X, name='x')
        s = model.binary_var_dict(keys=X, name='s')
    else:
        raise ValueError("Dominio non valido. Usare 'cont' o 'bin'.")
    return x, s
\end{lstlisting}
I vincoli principali del modello vengono aggiunti con la funzione \texttt{add\_constraints}, che implementa:
\begin{itemize}
    \item l'unicità di una configurazione attiva per ogni intervallo
    \item la compatibilità tra configurazioni consecutive
    \item le transizioni tra intervalli temporali
    \item i vincoli di permanenza e quiescenza
\end{itemize}
\clearpage

L'implementazione dei vincoli descritti è realizzata come segue:
\begin{lstlisting}[language=Python, caption={}]
def add_constraints(model, x, s, CONF_t, CONF_tc, T, tp, tq):
    # Vincolo: assegnazione di una configurazione per ogni istante
    for t in T:
        model.add_constraint(
            model.sum(x[(t, c)] for c in CONF_t[t]) == 1
        )

    # Vincolo: transizione tra configurazioni
    for t in T[:-1]:
        for c in CONF_t[t]:
            model.add_constraint(
                x[(t, c)] - model.sum(x[(t + 1, k)] for k in CONF_tc[(t + 1, c)]) <= 0
            )

    # Vincolo: cambiamento tra istanti consecutivi
    for t in T[1:]:
        for c in CONF_t[t]:
            if c in CONF_t[t - 1]:
                model.add_constraint(x[(t, c)] - x[(t - 1, c)] <= s[(t, c)])
            else:
                model.add_constraint(x[(t, c)] <= s[(t, c)])

    # Vincolo di permanenza
    if tp >= 2:
        for t in T:
            model.add_constraint(
                model.sum(
                    s[(w, c)] for w in range(t, min(t + tp, len(T))) 
                    for c in CONF_t[w]
                ) <= 1
            )

    # Vincolo di quiescenza
    if tq >= 2:
        for t in T:
            for c in CONF_t[t]:
                model.add_constraint(
                    x[(t, c)] + model.sum(
                        s[(w, c)] 
                        for w in range(t + 1, min(t + tq + 1, len(T))) 
                        if c in CONF_t[w]
                    ) <= 1)
\end{lstlisting}

Una volta assemblato il modello completo, la soluzione ottima viene ricercata invocando il metodo \texttt{model.solve()}. I risultati, che descrivono la pianificazione temporale delle configurazioni, vengono infine serializzati in formato JSON per le analisi successive.

\section{Verifica funzionale e analisi prestazionale}
A conclusione della fase implementativa, è stata condotta un'attività di verifica volta a validare il corretto funzionamento dell'intera pipeline software e la solvibilità del modello del \gls{dac}.
In assenza delle geometrie dei settori elementari, come discusso nella Sezione \ref{sec:limitiadrr}, la validazione si è basata su istanze sintetiche, costruite applicando le procedure descritte nella Sezione \ref{sec:gen}.

L'attività ha avuto il duplice obiettivo di confermare la correttezza logica del flusso di dati e di valutare l'efficienza computazionale dei tre stadi principali della pipeline: calcolo delle metriche, generazione delle configurazioni spaziali e risoluzione del problema di ottimizzazione.

Un primo risultato emerge dalla capacità del sistema di processare efficientemente grandi moli di dati. I dataset \gls{adrr} analizzati, relativi ad archi temporali mensili, raggiungono dimensioni considerevoli, contando fino a 7 milioni di record per singola mensilità. 
Grazie all'architettura basata su \texttt{Dask} e alla parallelizzazione delle operazioni di I/O, i tempi per il calcolo delle metriche di traffico (\emph{throughput} e \emph{occupancy}) e per la stima delle capacità sull'intero arco temporale analizzato (una mensilità) si sono attestati nell'ordine dei minuti ($\le 5$~minuti) su workstation standard. Questo risultato conferma che la scelta di evitare un \gls{dbms} relazionale in favore di un approccio orientato ai file (\emph{Parquet}) e al calcolo \emph{in-memory} è adeguata per trattare scenari su scala nazionale.

\subsection{Scalabilità della generazione degli scenari}
Anche la fase di costruzione dei pseudo-settori collassati, computazionalmente critica per via delle numerose operazioni geometriche (intersezioni e unioni di poligoni) e della natura combinatoria del problema, ha mostrato una buona scalabilità.
I test condotti hanno simulato la generazione di spazi delle configurazioni di diversa cardinalità. Anche negli scenari più onerosi i tempi di generazione non hanno superato l'ordine dei minuti ($\le 3$~minuti). Ciò dimostra l'efficacia dell'algoritmo di visita \gls{dfs} e dell'uso degli indici spaziali per abbattere la complessità delle query geometriche.

\subsection{Risoluzione del modello di ottimizzazione}
Infine, le istanze del problema di \gls{dac} così generate sono state risolte utilizzando il motore \gls{cplex}. I tempi di risoluzione registrati si sono mantenuti nell'ordine dei secondi, confermando che l'approccio \gls{ilp} rimane trattabile per istanze di queste dimensioni.

\subsection{Estensibilità e indipendenza dal dato geometrico}
La struttura del prototipo è stata progettata per essere indipendente dalla geometria specifica, rendendo la logica applicativa agnostica rispetto alla forma dei settori. Questa flessibilità permetterà l'impiego del software nel contesto reale non appena disponibili le geometrie ufficiali, senza richiedere modifiche sostanziali al codice. Ciò assicura un alto grado di riusabilità dello strumento, ponendo basi solide e scalabili per l'operatività futura.